\clearpage


\chapter{Технологическая часть}\label{ch:-2}

В этом разделе обоснованы средства реализации, а так же представлены реализации алгоритмов и функциональные тесты.


\section{Средства реализации}\label{sec:-2}

Для реализации алгоритмов в этой лабораторной работы был выбран язык $Python^{[1]}$ по следующим причинам:
\begin{itemize}
    \item python предоставляет необходимые структуры данных, такие как списки (массивы) и вложенные списки (матрицы);
    \item поддерживаемые библиотеки, такие как $matplotlib{[2]}$ позволяют визуализировать результаты исследования.
\end{itemize}

\clearpage


\section{Реализация алгоритмов}

В листингах~\ref{lst:std} ---~\ref{lst:optwyn} представлены реализации алгоритмов.

\begin{lstlisting}[style=python, label=lst:std,caption=Стандартный алгоритм умножения матриц]
def standard_multiply(a, b):
    rows_a = len(a)
    rows_b = len(b)

    if rows_a == 0 or rows_b == 0:
        return None

    cols_a = len(a[0])
    cols_b = len(b[0])

    if cols_a != rows_b or cols_a == 0 or cols_b == 0:
        return None

    c = [[0] * cols_b for _ in range(rows_a)]

    for i in range(rows_a):
        for j in range(cols_b):
            for k in range(cols_a):
                c[i][j] += a[i][k] * b[k][j]

    return c
\end{lstlisting}
\clearpage


\begin{lstlisting}[style=python, label=lst:wyn,caption=Алгоритм Винограда]
def winograd_multiply(a, b):
    rows_a = len(a)
    rows_b = len(b)

    if rows_a == 0 or rows_b == 0:
        return None

    cols_a = len(a[0])
    cols_b = len(b[0])

    if cols_a != rows_b or cols_a == 0 or cols_b == 0:
        return None

    c = [[0] * cols_b for _ in range(rows_a)]

    row_factor = [0] * rows_a
    col_factor = [0] * cols_b

    for i in range(rows_a):
        for j in range(cols_a // 2):
            row_factor[i] += a[i][2 * j] * a[i][2 * j + 1]

    for i in range(cols_b):
        for j in range(rows_b // 2):
            col_factor[i] += b[2 * j][i] * b[2 * j + 1][i]

    for i in range(rows_a):
        for j in range(cols_b):
            c[i][j] += -row_factor[i] - col_factor[j]

            for k in range(cols_a // 2):
                c[i][j] += (a[i][2 * k] + b[2 * k + 1][j]) * (a[i][2 * k + 1] + b[2 * k][j])

    if cols_a % 2 != 0:
        for i in range(rows_a):
            for j in range(cols_b):
                c[i][j] += a[i][cols_a - 1] * b[rows_b - 1][j]

    return c
\end{lstlisting}

\clearpage

\begin{lstlisting}[style=python, label=lst:optwyn,caption=Оптимизированный алгоритм Винограда]
def optimized_winograd_multiply(a, b):
    rows_a = len(a)
    rows_b = len(b)

    if rows_a == 0 or rows_b == 0:
        return None

    cols_a = len(a[0])
    cols_b = len(b[0])

    if cols_a != rows_b or cols_a == 0 or cols_b == 0:
        return None

    c = [[0] * cols_b for _ in range(rows_a)]

    row_factor = [0] * rows_a
    col_factor = [0] * cols_b

    for i in range(rows_a):
        j = cols_a // 2 - 1
        while j >= 0:
            row_factor[i] += a[i][j << 1] * a[i][(j << 1) + 1]
            j -= 1

    for i in range(cols_b):
        j = rows_b // 2 - 1
        while j >= 0:
            col_factor[i] += b[j << 1][i] * b[(j << 1) + 1][i]
            j -= 1

    for i in range(rows_a):
        for j in range(cols_b):
            c[i][j] += -row_factor[i] - col_factor[j]

            for k in range(cols_a // 2):
                c[i][j] += (a[i][k << 1] + b[(k << 1) + 1][j]) * (a[i][(k << 1) + 1] + b[k << 1][j])

            if cols_a % 2 != 0:
                c[i][j] += a[i][cols_a - 1] * b[rows_b - 1][j]

    return c
\end{lstlisting}

\clearpage


\section{Функциональные тесты}
В таблице~\ref{tbl:func_test} приведены функциональные тесты, на которых тестировалась программа.
\begin{table}[h]
    \begin{center}
        \caption{\label{tbl:func_test} Функциональные тесты, на которых были протестированы алгоритмы (в результате -- матрица c).}
        \begin{tabular}{|c|c|c|c|}
            \hline
            Матрица a & Матрица b & Результат
            \\ \hline
            $\begin{pmatrix}
                 1 & 2 \\
                 3 & 4
            \end{pmatrix}$ &
            $\begin{pmatrix}
                 5 & 6 \\
                 7 & 8
            \end{pmatrix}$ &
            $\begin{pmatrix}
                 19 & 22 \\
                 43 & 50
            \end{pmatrix}$
            \\ \hline
            $\begin{pmatrix}
                 1 & 2 & 3 \\
                 4 & 5 & 6
            \end{pmatrix}$ &
            $\begin{pmatrix}
                 7  & 8  \\
                 9  & 10 \\
                 11 & 12
            \end{pmatrix}$ &
            $\begin{pmatrix}
                 58  & 64  \\
                 139 & 154
            \end{pmatrix}$
            \\ \hline
            $\begin{pmatrix}
                 1 & 2 & 3
            \end{pmatrix}$ &
            $\begin{pmatrix}
                 4 \\
                 5 \\
                 6
            \end{pmatrix}$ &
            $\begin{pmatrix}
                 32
            \end{pmatrix}$ \\
            \hline
            $\begin{pmatrix}
                 1 & 2 & 3 \\
                 4 & 5 & 6 \\
                 7 & 8 & 9
            \end{pmatrix}$ &
            $\begin{pmatrix}
                 9 & 8 & 7 \\
                 6 & 5 & 4 \\
                 3 & 2 & 1
            \end{pmatrix}$ &
            $\begin{pmatrix}
                 30  & 24  & 18 \\
                 89  & 69  & 54 \\
                 138 & 114 & 90
            \end{pmatrix}$ \\
            \hline
            $\begin{pmatrix}
                 1 & 0 & 0 \\
                 0 & 1 & 0 \\
                 0 & 0 & 1
            \end{pmatrix}$ &
            $\begin{pmatrix}
                 5  & 6  & 7  \\
                 8  & 9  & 10 \\
                 11 & 12 & 13
            \end{pmatrix}$ &
            $\begin{pmatrix}
                 5  & 6  & 7  \\
                 8  & 9  & 10 \\
                 11 & 12 & 13
            \end{pmatrix}$ \\
            \hline
            $\begin{pmatrix}
                 \\
            \end{pmatrix}$ &
            $\begin{pmatrix}
                 \\
            \end{pmatrix}$ &
            \text{Ошибка} \\
            \hline
            $\begin{pmatrix}
                 1 & 2 \\
                 3 & 4
            \end{pmatrix}$ &
            $\begin{pmatrix}
                 5 & 6 & 7
            \end{pmatrix}$ &
            \text{Ошибка} \\
            \hline
        \end{tabular}
    \end{center}
\end{table}

\section*{Вывод}

В технологической части были представлены листинги реализованных алгоритмов, включая стандартный алгоритм умножения матриц,
а также алгоритм Винограда и оптимизированный алгоритм Винограда.
Также были проведены функциональные тесты, которые подтвердили корректность работы алгоритмов на различных входных данных, включая граничные случаи.
Результаты тестов показали, что разработанные реализации соответствуют поставленным требованиям.

